% !TEX program = xelatex
\documentclass[9pt,aspectratio=169]{beamer}
\usetheme{metropolis}

% ── Packages ──────────────────────────────────────────────
\usepackage{amsmath,amssymb}
\usepackage{booktabs}
\usepackage{graphicx}
\usepackage{tikz}
\usepackage{appendixnumberbeamer}

% ── Figure paths ──────────────────────────────────────────
\graphicspath{%
  {../output/diagnostics/figures/}%
  {../output/did/figures/}%
  {../output/bids/figures/}%
  {../notes/empirics/compra_agil_descriptives/}%
}

% ── Table path helper ─────────────────────────────────────
% All \input{} paths are relative to this file's directory

% ── Meta ──────────────────────────────────────────────────
\title{Local Preference Policies in Public Procurement:\\
       Efficiency, Distributional Effects, and Spillovers}
\subtitle{Evidence from Chile's Compra Ágil Reform}
\author{Nicolás Jiménez}
\date{March 2026}
\institute{Yale University}

\begin{document}

% ══════════════════════════════════════════════════════════
% TITLE
% ══════════════════════════════════════════════════════════
\maketitle

% ══════════════════════════════════════════════════════════
% INTRODUCTION
% ══════════════════════════════════════════════════════════
\section{Introduction}

% ── Motivation ───────────────────────────────────────────
\begin{frame}{Motivation}
\textbf{[PLACEHOLDER — fill with motivation slide]}

\medskip
Suggested content:
\begin{itemize}
  \item Public procurement is large: OECD governments spend 12--15\% of GDP
        on procurement. Developing countries often more.
  \item Local preference policies are widespread: Buy American, EU local content
        rules, developing-country SME preferences.
  \item Tension: local preferences may protect domestic firms and create local
        jobs, but can reduce competition and raise costs for the government.
  \item Little causal evidence on the competition and welfare effects of these
        policies, especially in developing countries.
  \item Chilean Compra Ágil reform (Dec 2024) provides a sharp, quasi-random
        change in the scope of local preference.
\end{itemize}
\end{frame}

% ── Research Questions ───────────────────────────────────
\begin{frame}{Research Questions}
\begin{enumerate}
  \item[\textbf{Q1}] \textbf{Efficiency:} What are the short-run effects of local
        preference policies on competition in public procurement?
        \begin{itemize}
          \item Number and composition of bidders
          \item Government procurement costs (winning bid prices)
          \item Heterogeneous effects by market thickness and distance
        \end{itemize}

  \vspace{0.5em}
  \item[\textbf{Q2}] \textbf{Distributional:} Who wins?
        \begin{itemize}
          \item Local vs.\ non-local firms; SME vs.\ large firms
          \item Do local entry gains offset non-local exclusion?
        \end{itemize}

  \vspace{0.5em}
  \item[\textbf{Q3}] \textbf{Spillovers:} Do local preference policies in
        peripheral markets affect competition in urban markets?
        \begin{itemize}
          \item Capacity reallocation: freed firms re-enter home markets
          \item Cross-market effects missed by single-market analyses
        \end{itemize}
\end{enumerate}
\end{frame}

% ── Project Overview ─────────────────────────────────────
\begin{frame}{This Project}
\begin{columns}[T]
  \column{0.5\textwidth}
  \textbf{Part I: Empirics}
  \begin{itemize}
    \item Descriptive facts about local and non-local firm participation
          in Chilean public procurement
    \item DiD estimates of the short-run effects of the Compra Ágil reform
          on bidder entry, bidder composition, and winning prices
    \item Regional heterogeneity: how effects vary with pre-reform market
          structure and geography
    \item Cross-market spillovers: competition effects in \textit{above}-threshold
          markets driven by firm capacity reallocation
  \end{itemize}

  \column{0.5\textwidth}
  \textbf{Part II: Structural Model}
  \begin{itemize}
    \item Spatial model of firm entry and bidding for public procurement projects
    \item Firms have home regions; face distance-based cost and entry cost
          disadvantages; are subject to capacity constraints linking markets
    \item Estimation using pre- and post-reform variation
    \item \textbf{Counterfactuals:}
    \begin{itemize}
      \item[-] Optimal local preference threshold by region
      \item[-] Welfare decomposition (cost vs.\ local benefit)
      \item[-] Alternative policy designs
    \end{itemize}
  \end{itemize}
\end{columns}
\end{frame}

% ── Related Literature ───────────────────────────────────
\begin{frame}{Related Literature}
\textbf{[PLACEHOLDER — fill with related literature]}

\medskip
Key strands:
\begin{itemize}
  \item \textbf{Public procurement and competition:}
        Athey, Levin \& Seira (2011); Krasnokutskaya \& Seim (2011);
        Decarolis (2014); Coviello \& Mariniello (2014)
  \item \textbf{Local preference / Buy-national policies:}
        Marion (2007); Kang \& Miller (2022); Shingal (2020)
  \item \textbf{SME preferences in procurement:}
        Carril, Gonzalez-Lira \& Walker (2022); Sorensen (2022)
  \item \textbf{Procurement reform in developing countries:}
        Bandiera et al.\ (2009); Best et al.\ (2017); Bosio et al.\ (2022)
  \item \textbf{Cross-market spillovers and capacity constraints:}
        Jeziorski \& Krasnokutskaya (2016); Groeger (2014)
\end{itemize}

\medskip
\textbf{Contribution:} First quasi-experimental estimates of local preference
effects accounting for cross-market spillovers via capacity reallocation.
Structural model quantifies welfare trade-offs and enables counterfactuals.
\end{frame}


% ══════════════════════════════════════════════════════════
% POLICY AND SETTING
% ══════════════════════════════════════════════════════════
\section{Setting}

% ── Reform Overview ──────────────────────────────────────
\begin{frame}{Compra Ágil Reform (Ley 21.634)}
\begin{columns}[T]
  \column{0.49\textwidth}
  \textbf{Before (until Dec 11, 2024)}
  \begin{itemize}
    \item Compra Ágil limit: \textbf{30 UTM} ($\approx$ CLP 2M, USD 2k)
    \item Projects 30--100 UTM went through competitive licitaciones
    \item No explicit SME/local priority at first stage
    \item Non-local firms competed freely in all markets
  \end{itemize}
  \column{0.49\textwidth}
  \textbf{After (from Dec 12, 2024)}
  \begin{itemize}
    \item Compra Ágil limit: \textbf{100 UTM} ($\approx$ CLP 7M, USD 7k)
    \item Projects 30--100 UTM now use simplified process
    \item \textbf{SME and local firm priority} in first instance
    \item Non-local firms effectively excluded from this band
  \end{itemize}
\end{columns}

\bigskip
\textbf{Timeline:} Law 21.634 published November 2024; effective December 12, 2024.
First post-reform month with full exposure: January 2025.

\bigskip
\textbf{Scale:} The 30--100 UTM band accounts for $\approx$\,25--35\% of all
municipal procurement tenders by volume. The reform substantially narrows the
set of markets where non-local firms can participate.
\end{frame}

% ── Reform Impact on Procurement Mode ────────────────────
\begin{frame}{Projects Shifted Almost Entirely to Compra Ágil}
\centering
\includegraphics[width=0.96\linewidth]{ca_D4_share_compra_agil_by_utm_bucket.png}

\vspace{0.5em}
\small\textbf{Takeaway:} Nearly all 30--100 UTM projects switched from
licitaciones to Compra Ágil after the reform --- compliance is near-complete.
\end{frame}

% ── Data ─────────────────────────────────────────────────
\begin{frame}{Data}
\begin{columns}[T]
  \column{0.5\textwidth}
  \textbf{ChileCompra (Mercado Público)}
  \begin{itemize}
    \item Universe of public procurement contracts in Chile
    \item Tender-level: estimated value (UTM), procurement type,
          region, sector, date
    \item Bid-level: firm RUT, bid amount, winner indicator
    \item \textbf{Coverage:} Jan 2022 -- Mar 2025
    \item \textbf{Sample:} Municipalidades sector;
          treated band (30--100 UTM);
          control bands (0--30, 100--200 UTM)
  \end{itemize}

  \column{0.5\textwidth}
  \textbf{Servicio de Impuestos Internos (SII)}
  \begin{itemize}
    \item Chilean Internal Revenue Service administrative data
    \item Firm-level: size classification (micro, small, medium, large),
          sector, location, annual sales
    \item \textbf{Use:} Classify bidders as SME (MiPYME) vs.\ large;
          local vs.\ non-local (matched to buyer region)
    \item \textbf{Coverage:} Annual data, matched to ChileCompra via RUT
  \end{itemize}
\end{columns}

\bigskip
\textbf{Final sample (Municipalidades):} $\approx$\,2,500 procuring entities;
$\approx$\,180,000 tenders; $\approx$\,650,000 bids over the study period.
\end{frame}

% ── Summary Statistics: Procurement ──────────────────────
\begin{frame}{Sample Overview: Pre-Reform Descriptives}
\textbf{[PLACEHOLDER — fill with summary statistics table]}

\medskip
Suggested content:
\begin{itemize}
  \item Table: N tenders, N bidders (mean), SME win share, log price ratio
        --- split by sector (Municipalidades vs.\ Obras Públicas) and UTM band
        (0--30, 30--100, 100--200 UTM), pre-reform only.
  \item Highlight that the treated band (30--100 UTM) had systematically
        lower bidder counts and higher non-local shares than control bands.
  \item Source: \texttt{did\_tender\_sample.parquet} + \texttt{did\_bid\_sample.parquet}.
\end{itemize}
\end{frame}

% ── Summary Statistics: SII ──────────────────────────────
\begin{frame}{Sample Overview: Bidder Characteristics (SII Data)}
\textbf{[PLACEHOLDER — fill with SII summary statistics]}

\medskip
Suggested content:
\begin{itemize}
  \item Share of bids by SME vs.\ large firm, by UTM band and sector
  \item Share of bids by local vs.\ non-local firm (home region = buyer region)
  \item Average firm size (sales, employees) of winning vs.\ losing bidders
  \item Geographic concentration: fraction of bidders from RM bidding in
        other regions (``export share'')
\end{itemize}
\end{frame}


% ══════════════════════════════════════════════════════════
% EMPIRICS
% ══════════════════════════════════════════════════════════
\section{Empirical Strategy \& Results}

% ── Facts About Local/Non-local Firms ────────────────────
\begin{frame}{Descriptive Facts: Local vs.\ Non-Local Bidders}
\textbf{[PLACEHOLDER --- fill with descriptive facts about local and non-local firm bids]}

\medskip
Suggested content:
\begin{itemize}
  \item Bar/line chart: share of bids by local vs.\ non-local firms by
        UTM band and region, pre-reform
  \item Map or bar chart: which regions send firms to other regions most?
        (foreshadows the export share analysis)
  \item Heterogeneity: non-local share is higher in remote regions,
        lower in RM (the inverse-U pattern from the binscatter)
\end{itemize}
\end{frame}

% ── DiD Specification ────────────────────────────────────
\begin{frame}{Difference-in-Differences: Specification}
\textbf{Estimating equation (TWFE):}
\[
  y_{it} = \alpha_i + \gamma_t + \beta\,\mathbf{1}[\text{treated}_i \times \text{post}_t] + \varepsilon_{it}
\]

\medskip
\textbf{Treatment and control groups} (defined by estimated contract value in UTM):
\begin{columns}[T]
  \column{0.48\textwidth}
  \begin{itemize}
    \item \textbf{Treated:} 30--100 UTM \\
          \small Tenders that switch from competitive bidding (licitación)
          to Compra Ágil after Dec 12, 2024.
  \end{itemize}
  \column{0.52\textwidth}
  \begin{itemize}
    \item \textbf{Control 1} (low): 1--30 UTM \\
          \small Already under Compra Ágil before the reform.
    \item \textbf{Control 2} (high): 100--200 UTM \\
          \small Above the new threshold; procurement process unchanged.
  \end{itemize}
\end{columns}

\medskip
\textbf{Fixed effects:} Procuring entity $\alpha_i$ + year-month $\gamma_t$.\\
\textbf{Standard errors:} Clustered by procuring entity.\\
\textbf{Sample:} Jan 2023 -- Mar 2025; Municipalidades sector.\\[4pt]
\small Results show control group: \textit{Control 2 (100--200 UTM)}.
\end{frame}

% ── IV-DiD Specification ─────────────────────────────────
\begin{frame}{IV-DiD: Specification}
\textbf{Motivation:} Intent-to-treat (DiD) estimates may under-state effects if
compliance is imperfect. We instrument actual Compra Ágil use with reform
eligibility.

\bigskip
\textbf{First stage (within-entity, within-month):}
\[
  \widetilde{CA}_{it} = \pi\,\widetilde{Z}_{it} + \nu_{it},
  \qquad
  Z_{it} = \mathbf{1}[\text{treated}_i \times \text{post}_t]
\]

\textbf{Second stage:}
\[
  \tilde{y}_{it} = \beta_{\text{IV}}\,\widehat{\widetilde{CA}}_{it} + \varepsilon_{it}
\]

\begin{itemize}
  \item \textbf{Endogenous:} $CA_{it}$ --- indicator, tender processed via Compra Ágil
  \item \textbf{Instrument:} $Z_{it} = \text{Treated}_i \times \text{Post}_t$ --- reform eligibility
  \item Same FE and clustering as OLS DiD; Control 2 only (100--200 UTM)
  \item $\beta_{\text{IV}}$ is a LATE: effect of Compra Ágil \emph{use}
\end{itemize}
\end{frame}

% ── IV First Stage ───────────────────────────────────────
\begin{frame}{IV-DiD: First Stage (Municipalidades)}
\small
\medskip
\input{../output/did/tables/did_table_first_stage_munic.tex}
\end{frame}

% ── Event Studies ─────────────────────────────────────────
\begin{frame}{Event Studies: N Bidders and SME Share (Municipalidades)}
\begin{columns}[c]
  \column{0.5\textwidth}\centering
  \includegraphics[width=\linewidth]{../output/did/figures/diag_event_study_n_bidders_munic.png}
  \column{0.5\textwidth}\centering
  \includegraphics[width=\linewidth]{../output/did/figures/diag_event_study_sme_share_sii_munic.png}
\end{columns}
\end{frame}

\begin{frame}{Event Studies: SME Winners and Win Price (Municipalidades)}
\begin{columns}[c]
  \column{0.5\textwidth}\centering
  \includegraphics[width=\linewidth]{../output/did/figures/diag_event_study_winner_is_sme_sii_munic.png}
  \column{0.5\textwidth}\centering
  \includegraphics[width=\linewidth]{../output/did/figures/diag_event_study_log_win_price_ratio_munic.png}
\end{columns}
\end{frame}

\begin{frame}{Event Study: Bid-Level Price (Municipalidades)}
\centering
\includegraphics[width=0.58\linewidth]{../output/did/figures/diag_event_study_log_sub_price_ratio_munic.png}

\medskip
\small\textit{Flat pre-trends; significant post-reform increase in winning and
submitted prices --- consistent with reduced competition (fewer non-local firms)
not fully offset by local entry.}
\end{frame}

% ── IV-DiD Results ───────────────────────────────────────
\begin{frame}{IV-DiD: Entry at Bidding Stage (Panel A, Municipalidades)}
\small\medskip
\input{../output/did/tables/did_table_compare_iv_A_munic.tex}
\end{frame}

\begin{frame}{IV-DiD: Bidder Composition (Panel B, Municipalidades)}
\small\medskip
\input{../output/did/tables/did_table_compare_iv_B_munic.tex}
\end{frame}

\begin{frame}{IV-DiD: Winner Characteristics (Panel C, Municipalidades)}
\small\medskip
\input{../output/did/tables/did_table_compare_iv_C_munic.tex}
\end{frame}

\begin{frame}{IV-DiD: Bid Outcomes (Panel D, Municipalidades)}
\small\medskip
\input{../output/did/tables/did_table_compare_iv_D_munic.tex}
\end{frame}

% ── Bid markups ──────────────────────────────────────────
\begin{frame}{Bid Markups: Within-Auction Comparison}
\centering
\includegraphics[width=0.62\linewidth]{bids_part1_coefplot.png}

\medskip
\scriptsize
\begin{itemize}
  \item Conditional on \textbf{auction fixed effects}, local and non-local firms submit \textbf{very similar markups}; the local coefficient is near zero.
  \item This is a \textbf{conditional-on-entry} comparison: among firms that actually bid in the same auction, non-local firms do not look much more expensive.
  \item The main cross-sectional difference is by \textbf{firm size}: SMEs bid about $0.26$ log points higher relative to estimated cost.
\end{itemize}
\end{frame}

\begin{frame}{Bid Markups: Within-Firm Home-Market Advantage}
\centering
\includegraphics[width=0.62\linewidth]{bids_part2_coefplot.png}

\medskip
\scriptsize
\begin{itemize}
  \item For the \textbf{same firm}, local bids are lower: about $-0.086$ log points with bidder and month fixed effects.
  \item With buyer-region fixed effects added, the local coefficient remains negative at about $-0.040$, so the pattern is not just region composition.
  \item The \textbf{distance slope is estimated separately on non-local bids only}: $\log(\text{distance})$ enters without the local dummy, and its coefficient is about $-0.038$.
  \item Interpretation: the same firm bids differently at home vs.\ away, so there is a genuine \textbf{home-market advantage / non-local friction}.
\end{itemize}
\end{frame}

\begin{frame}{Bid Markups: Within-Firm Differences by Firm Size}
\centering
\includegraphics[width=0.68\linewidth]{bids_part2_size_split_local.png}

\medskip
\scriptsize
\begin{itemize}
  \item Both \textbf{SME} and \textbf{large} firms bid less on local projects when we compare the same firm across auctions.
  \item The point estimate is \textbf{more negative for large firms} ($-0.126$ vs.\ $-0.076$ with bidder + month FE), so we do \textbf{not} see evidence that large firms have a smaller home-market wedge.
  \item The within-firm local advantage is therefore \textbf{not only an SME phenomenon}; it appears across firm types.
\end{itemize}
\end{frame}

\begin{frame}{Bid Markups: Interpretation for the Model}
\small
\textbf{What the bid-markup evidence says:}
\begin{itemize}
  \item \textbf{Within firm:} the same supplier bids lower at home. This identifies a \textbf{local advantage} or, equivalently, a \textbf{non-local cost/information wedge}.
  \item \textbf{Within auction:} local and non-local bidders look similar conditional on entry. Observed non-local entrants must therefore be a \textbf{positively selected} set.
  \item The natural model combines:
        \begin{itemize}
          \item firm heterogeneity / efficiency,
          \item a distance or home-market friction,
          \item endogenous entry into non-local auctions.
        \end{itemize}
  \item So the reduced-form evidence is consistent with \textbf{real non-local frictions} that are partly offset in cross-section by the selection of stronger firms into away markets.
\end{itemize}
\end{frame}

\begin{frame}{Bid Markups: Reform Effect in Municipal Bids}
\centering
\includegraphics[width=0.62\linewidth]{bids_part3_did_coefplot.png}

\medskip
\scriptsize
\begin{itemize}
  \item In the treated 30--100 UTM band, bid markups rise by about \textbf{$0.40$ log points} after the reform relative to 100--200 UTM auctions.
  \item The \textbf{local triple interaction is small and insignificant}: the markup increase is not concentrated only among local bids.
  \item The \textbf{SME triple interaction is positive}, suggesting the reform-era markup increase may be somewhat larger among SME bidders.
\end{itemize}
\end{frame}

% ── Validity Checks ───────────────────────────────────────
\begin{frame}{DiD Validity Checks}
\textbf{Three diagnostic tests (Control 2 = 100--200 UTM):}
\begin{enumerate}
  \item \textbf{Wald pre-trend F-test} --- joint test
        $H_0: \beta_{-8} = \cdots = \beta_{-2} = 0$
        using full cluster-robust covariance matrix of event-study coefficients
  \item \textbf{Time placebo} --- DiD run on pre-reform data only
        ($k \in [-8,-1]$), fake reform at $k = -4$; coefficient should be
        $\approx 0$
  \item \textbf{Pre-reform balance} --- OLS of entity-level pre-reform means
        on Treated; tests level differences before reform
\end{enumerate}
\end{frame}

\begin{frame}{DiD: Pre-Trend Wald Test (Municipalidades)}
\small\medskip
\input{../output/did/tables/diag_pretrend_wald_munic.tex}
\end{frame}

\begin{frame}{DiD: Time Placebo Test (Municipalidades)}
\small\medskip
\input{../output/did/tables/diag_placebo_munic.tex}
\end{frame}

\begin{frame}{DiD: Pre-Reform Balance (Municipalidades)}
\small\medskip
\input{../output/did/tables/diag_balance_munic.tex}
\end{frame}


% ══════════════════════════════════════════════════════════
% REGIONAL HETEROGENEITY
% ══════════════════════════════════════════════════════════
\section{Regional Heterogeneity}

\begin{frame}{Regional Heterogeneity: Market Structure Varies by Region}
\begin{columns}[c]
  \column{0.52\textwidth}
  \centering
  \includegraphics[width=\linewidth]{binscatter_q_nonlocal_munic.png}
  \column{0.48\textwidth}
  \textbf{Four pre-reform moderators (Municipalidades, 30--100 UTM):}
  \begin{itemize}
    \item \textbf{Nonlocal share:} fraction of bids by out-of-region firms
    \item \textbf{q\_pre:} avg.\ monthly tender count
    \item \textbf{n\_pot\_local:} unique local suppliers
    \item \textbf{Distance from Santiago:} geographic isolation (km)
  \end{itemize}
  \medskip
  \textit{Thick markets (high q\_pre) have lower non-local penetration ---
  market size proxies for local competitive capacity.}
\end{columns}
\end{frame}

\begin{frame}{Regional Heterogeneity: Specification}
\textbf{Question:} Does the effect of Compra Ágil vary with pre-reform
regional characteristics?

\bigskip
\textbf{Estimating equation (OLS DiD):}
\[
  y_{irt} = \alpha_i + \gamma_t
    + \beta_1\,(T_i \times \text{Post}_t)
    + \beta_2\,(T_i \times \text{Post}_t \times Z_r)
    + \varepsilon_{irt}
\]

\medskip
\begin{itemize}
  \item $Z_r$ = standardised pre-reform regional moderator
        (one specification per moderator)
  \item $\beta_1$ = average treatment effect;
        \quad $\beta_2$ = heterogeneity w.r.t.\ $Z_r$
  \item Same FE and clustering as baseline DiD
        (entity + year-month FE; SEs by PE)
  \item \textbf{Sample:} Municipalidades sector
\end{itemize}

\medskip
\textbf{Five moderators $Z_r$:} non-local share, tender count, tender value,
potential local firms, distance from Santiago.
\end{frame}

\begin{frame}{Exposure-Based Design: Panel Version of High vs.\ Low Exposure}
\small
\textbf{Idea:} compare places that were more exposed to non-local competition
\emph{before} the reform to places that were less exposed.

\medskip
\textbf{Why use a panel instead of post-only comparisons?}
\begin{itemize}
  \item Exposure is measured \textbf{pre-reform}, so it is predetermined.
  \item Entity and month fixed effects absorb permanent level differences across places and common national shocks.
  \item The key coefficient is the \textbf{triple interaction}:
\[
  \beta_2 \;:\; \text{Treated} \times \text{Post} \times \text{Exposure}_r
\]
  \item Interpretation: do treated-band outcomes move \textbf{more} after the reform in high-exposure places?
\end{itemize}

\medskip
\textbf{Two exposure measures shown next:}
\begin{itemize}
  \item Direct measure: pre-reform non-local bidder share
  \item Geographic proxy: distance from Santiago
\end{itemize}
\end{frame}

\begin{frame}{Exposure-Based DiD: Pre-Reform Non-Local Share}
\scriptsize
\resizebox{\textwidth}{!}{%
\input{../output/did/tables/hetero_interacted_nonlocal_share_pre_munic.tex}}

\medskip
\small\textit{This is the table version of the exposure design. $\hat\beta_2$ is the additional reform effect in regions with higher pre-reform non-local bidder share. The clearest heterogeneity appears in bidder composition, lower single-bidder risk, and higher submitted bid markups in high-exposure regions.}
\end{frame}

\begin{frame}{Exposure-Based DiD: Distance from Santiago}
\scriptsize
\resizebox{\textwidth}{!}{%
\input{../output/did/tables/hetero_interacted_dist_from_santiago_munic.tex}}

\medskip
\small\textit{Distance from Santiago is a predetermined geographic proxy for exposure to non-local competition. The interacted coefficient $\hat\beta_2$ mainly loads on bidder composition and single-bidder probability; price heterogeneity is weaker than with the direct non-local-share measure.}
\end{frame}

\begin{frame}{Heterogeneity: N Local and N Nonlocal Bidders}
\begin{columns}[c]
  \column{0.5\textwidth}\centering
  \textbf{\small N Local Bidders}\\[-2pt]
  \includegraphics[width=0.9\linewidth]{hetero_coefplot_n_local_nonlocal_share_pre_munic.png}
  \column{0.5\textwidth}\centering
  \textbf{\small N Non-local Bidders}\\[-2pt]
  \includegraphics[width=0.9\linewidth]{hetero_coefplot_n_nonlocal_nonlocal_share_pre_munic.png}
\end{columns}
\medskip
\small\textit{Regions with high pre-reform non-local share see the largest local
entry gains --- consistent with the exclusion mechanism.}
\end{frame}

\begin{frame}{Heterogeneity: Total Bidders and SME Share (by Distance)}
\begin{columns}[c]
  \column{0.5\textwidth}\centering
  \textbf{\small N Bidders}\\[-2pt]
  \includegraphics[width=0.9\linewidth]{hetero_coefplot_n_bidders_dist_from_santiago_munic.png}
  \column{0.5\textwidth}\centering
  \textbf{\small SME Share (SII)}\\[-2pt]
  \includegraphics[width=0.9\linewidth]{hetero_coefplot_sme_share_sii_dist_from_santiago_munic.png}
\end{columns}
\medskip
\small\textit{Distance from Santiago is an exogenous geographic instrument for
pre-reform non-local penetration.}
\end{frame}

\begin{frame}{Heterogeneity: Win Price Ratio}
\begin{columns}[c]
  \column{0.5\textwidth}\centering
  \textbf{\small by Nonlocal Share}\\[-2pt]
  \includegraphics[width=0.9\linewidth]{hetero_coefplot_log_win_price_ratio_nonlocal_share_pre_munic.png}
  \column{0.5\textwidth}\centering
  \textbf{\small by Distance from Santiago}\\[-2pt]
  \includegraphics[width=0.9\linewidth]{hetero_coefplot_log_win_price_ratio_dist_from_santiago_munic.png}
\end{columns}
\medskip
\small\textit{Win prices rise most in regions that had high pre-reform
non-local competition --- local entry cannot fully replace lower-cost
non-local incumbents.}
\end{frame}


% ══════════════════════════════════════════════════════════
% CROSS-MARKET SPILLOVERS
% ══════════════════════════════════════════════════════════
\section{Cross-Market Spillovers}

\begin{frame}{Cross-Market Spillovers: Mechanism}
\textbf{Research question:} Did firms freed from 30--100 UTM licitaciones
redirect capacity to \textit{above-threshold} (100--500 UTM) tenders in their
home regions, increasing competition there?

\bigskip
\textbf{Mechanism (spatial auction model):}
\begin{enumerate}
  \item Reform excludes non-local firms from 30--100 UTM Compra Ágil.
  \item Firms from ``export regions'' (e.g.\ RM) had many cross-region bids
        in that band before the reform --- these bids are now blocked.
  \item Freed capacity $\Rightarrow$ these firms bid more in their home
        region's above-threshold (100--500 UTM) tenders.
  \item Prediction: N bidders (esp.\ large firms) $\uparrow$ post-reform
        in high-export-share regions.
\end{enumerate}

\bigskip
\textbf{Treatment:} \textit{Export share} =
fraction of pre-reform bids by home-region-$r$ firms placed in \textit{other}
regions (30--100 UTM band). Time-invariant; variation is cross-regional.
\end{frame}

\begin{frame}{Export Share by Region (Pre-Reform)}
\centering
\includegraphics[width=0.70\linewidth]{spillover_export_share_munic.png}

\medskip
\small\textit{Región Metropolitana has the highest export share: Santiago firms
dominated peripheral-region procurement in the treated band.
Far south and north regions are largely self-contained.}
\end{frame}

\begin{frame}{Cross-Market Spillovers: Specification}
\textbf{Event-study TWFE} on above-threshold (100--500 UTM) region--month panel:
\[
  y_{rt} = \alpha_r + \gamma_t
    + \sum_{k \neq -1} \delta_k \,\bigl(\mathbf{1}[t{=}k] \times
      \mathrm{ExportShare}_r\bigr) + \varepsilon_{rt}
\]

\begin{itemize}
  \item $r$ = region (16 Chilean regions), $t$ = year-month.
  \item $k$ = months relative to reform ($k=0$ = December 2024).
  \item $\alpha_r$ = region FE; $\gamma_t$ = year-month FE.
  \item Estimated for \textbf{continuous} (standardized) and
        \textbf{binary} (top-tercile) export share.
  \item Outcomes: N bidders, N local, N non-local, N large, N SME, Pr(single).
  \item SE clustered at region level (16 clusters).
\end{itemize}
\end{frame}

\begin{frame}{Spillovers: N Bidders and Pr(Single Bidder)}
\begin{columns}[c]
  \column{0.5\textwidth}\centering
  \includegraphics[width=\linewidth]{spillover_event_study_n_bidders_munic.png}
  \column{0.5\textwidth}\centering
  \includegraphics[width=\linewidth]{spillover_event_study_single_bidder_munic.png}
\end{columns}
\medskip
\small\textit{Total competition rises in high-export regions post-reform;
single-bidder probability falls --- consistent with freed-capacity reallocation.}
\end{frame}

\begin{frame}{Spillovers: Local vs.\ Non-Local Bidders}
\begin{columns}[c]
  \column{0.5\textwidth}\centering
  \includegraphics[width=\linewidth]{spillover_event_study_n_local_munic.png}
  \column{0.5\textwidth}\centering
  \includegraphics[width=\linewidth]{spillover_event_study_n_nonlocal_munic.png}
\end{columns}
\medskip
\small\textit{Local bidder count rises in high-export regions (home firms
re-enter above-threshold band). Non-local count adjusts as
previously-exporting firms redirect to home markets.}
\end{frame}

\begin{frame}{Spillovers: Large Firms and Non-Local Share}
\begin{columns}[c]
  \column{0.5\textwidth}\centering
  \includegraphics[width=\linewidth]{spillover_event_study_n_large_munic.png}
  \column{0.5\textwidth}\centering
  \includegraphics[width=\linewidth]{spillover_event_study_share_nonlocal_munic.png}
\end{columns}
\medskip
\small\textit{Large-firm entry is the main driver of the spillover effect.
Non-local share of above-threshold tenders rises in export regions as
home-based large firms redirect bids.}
\end{frame}

\begin{frame}{Spillovers: SME Bidders}
\centering
\includegraphics[width=0.60\linewidth]{spillover_event_study_n_sme_munic.png}

\medskip
\small\textit{SME spillover effect is smaller --- SMEs tend to be local and
were not the primary ``exporting'' firms. Effect concentrated among large firms.}
\end{frame}


% ══════════════════════════════════════════════════════════
% STRUCTURAL MODEL
% ══════════════════════════════════════════════════════════
\section{Structural Model}

% ── Environment ──────────────────────────────────────────
\begin{frame}{Model: Environment and Timing}
\begin{columns}[T]
  \column{0.5\textwidth}
  \textbf{Regions and firms}
  \begin{itemize}
    \item $R$ regions with pairwise distances $d_{rs}$
    \item Each period, region $r$ posts $Q_r$ procurement projects
          at reserve price $\bar{p}$
    \item Firms $i \in \mathcal{F}$ have home region $h_i$ and
          capacity $K_i$ (max markets per period)
    \item $N_r^{pot}$: potential local suppliers in region $r$
  \end{itemize}

  \column{0.5\textwidth}
  \textbf{Timing}
  \begin{enumerate}
    \item \textbf{Demand} $\mathbf{Q}$ realized and observed
    \item \textbf{Entry:} firm $i$ chooses set $M_i$,
          $|M_i| \leq K_i$; pays entry cost $\kappa_r^i$
    \item \textbf{Cost realization:} private cost drawn
          for each entered market
    \item \textbf{Bidding:} sealed-bid first-price auction;
          lowest bidder wins
  \end{enumerate}
\end{columns}
\end{frame}

% ── Cost Structure ───────────────────────────────────────
\begin{frame}{Model: Cost Structure}
\textbf{Firm $i$'s cost for a project in region $r$:}
\[
  c_{ir} = \underbrace{\bar{c}}_{\text{baseline}}
         + \underbrace{\delta \cdot d_{h_i, r}}_{\text{distance cost}}
         - \underbrace{\lambda \cdot \mathbf{1}[h_i = r]}_{\text{local advantage}}
         + \underbrace{\epsilon_{ir}}_{\text{private shock}}
\]

\medskip
\begin{itemize}
  \item $\bar{c} > 0$: common baseline project cost
  \item $\delta > 0$: per-unit-distance cost (transport, logistics,
        remote supervision)
  \item $\lambda > 0$: local informational and relational advantage
        (knowledge of local regulations, subcontractor networks,
        permitting relationships)
  \item $\epsilon_{ir} \sim F(\cdot)$ i.i.d.: idiosyncratic private cost shock
\end{itemize}

\medskip
\textbf{Key:} A non-local firm from region $s \neq r$ faces an expected cost
disadvantage of $\delta \cdot d_{sr} + \lambda$ relative to a local firm.
Distance cost and informational advantage are \emph{separately} identified.
\end{frame}

% ── Entry Costs ──────────────────────────────────────────
\begin{frame}{Model: Entry Costs and Capacity}
\textbf{Entry cost for firm $i$ into market $r$:}
\[
  \kappa_r^i = \kappa_0 + \kappa_d \cdot d_{h_i, r}
\]

\begin{itemize}
  \item $\kappa_0 > 0$: baseline participation cost (preparing bids,
        learning about projects) --- same for all firms
  \item $\kappa_d > 0$: additional cost of distant operations
  \item Local firms pay only $\kappa_0$
\end{itemize}

\bigskip
\textbf{Capacity constraint ($K_i$):}

Firm $i$ can enter at most $K_i$ markets per period. This is the key ingredient
that \textit{links markets}: a firm bidding in peripheral regions uses up
capacity that could be deployed at home.

\medskip
Entry set choice:
\[
  \max_{M_i \subseteq \{1,\dots,R\},\; |M_i| \leq K_i}
    \sum_{r \in M_i} \pi_{ir}(n_r^e)
\]
\end{frame}

% ── Policy Environment ───────────────────────────────────
\begin{frame}{Model: Policy Environment}
\textbf{Pre-reform:} All firms may enter any market.

\bigskip
\textbf{Post-reform --- three tiers by project value:}

\medskip
\begin{enumerate}
  \item \textbf{Strong preference (Compra Ágil, $< \underline{v}$ = 100 UTM):}\\
        Only local firms + EMT may enter. Non-local firms \textbf{fully excluded}.

  \vspace{0.5em}
  \item \textbf{Weak preference ($\underline{v}$--$\bar{v}$, 100--500 UTM):}\\
        All firms may enter, but local firms receive bid preference $\alpha > 0$.\\
        Non-local firm $i$ wins only if $b_i < b_j - \alpha$ for all locals $j$.

  \vspace{0.5em}
  \item \textbf{No preference ($> \bar{v}$ = 500 UTM):}\\
        Status quo; all firms compete on equal terms.
\end{enumerate}

\bigskip
\small\textbf{Note:} Theoretical analysis focuses on the strong-form case
(full exclusion, 30--100 UTM band) --- this is the empirically relevant tier
and generates the starkest predictions.
\end{frame}

% ── Equilibrium ─────────────────────────────────────────
\begin{frame}{Model: Equilibrium}
\textbf{Bidding stage (conditional on $n_r$ entrants):}
\[
  E[p_r \mid n]
    = \underbrace{E[c_{(1:n)}]}_{\text{efficient cost}}
    + \underbrace{\frac{1}{n} E\!\left[\frac{F(c_{(1:n)})}{f(c_{(1:n)})}\right]}_{\text{markup (decreasing in }n)}
\]
More entrants $\Rightarrow$ lower prices via both lower $c_{(1:n)}$ and lower markup.

\bigskip
\textbf{Entry equilibrium:} vector $(n_1^*, \dots, n_R^*)$ such that
for each market $r$ and each entering firm $i$:
\[
  Q_r \cdot \Pi_{ir}^{bid}(n_r^*) \geq \kappa_r^i
  \quad \text{(with equality for marginal entrant)}
\]

\medskip
\textbf{Multi-market:} the capacity constraint means
\[
  n_r^*\!\left(\{n_s^*\}_{s \neq r}\right) \quad \forall\; r
\]
must be solved as a system --- markets are interdependent through firm capacity.
\end{frame}

% ── Policy Effects: Direct ───────────────────────────────
\begin{frame}{Model: Direct Effect of the Policy}
\textbf{Immediate effect:} bidding pool in market $r$ shrinks from
$n_r^{pre}$ to $n_r^L$ (local entrants only).

\bigskip
Government pays more if excluded non-local firms were cost-competitive.
How much competition recovers depends on three factors:

\medskip
\begin{enumerate}
  \item \textbf{$N_r^{pot}$ (local firm pool):} if large relative to equilibrium
        entry, new entry can fully replace excluded firms. If small (remote regions),
        it cannot.
  \vspace{0.3em}
  \item \textbf{$Q_r$ (demand volume):} thick demand sustains more entrants.
        Thin, lumpy demand cannot support competitive local markets.
  \vspace{0.3em}
  \item \textbf{$\delta \cdot d_{sr} + \lambda$ (cost gap):} if excluded firms
        were from nearby regions, the government loses genuinely competitive bidders.
        If from far away, they were high-cost anyway.
\end{enumerate}
\end{frame}

% ── Cross-Market Spillover ────────────────────────────────
\begin{frame}{Model: Cross-Market Spillovers via Capacity}
\textbf{Mechanism:} Consider Santiago-based firm $i$ that pre-reform entered
both Santiago and rural region $r$ (using 2 of $K$ capacity units). Post-reform:

\begin{enumerate}
  \item Firm $i$ is excluded from $r$ --- recovers one capacity unit.
  \item If there is another market $s$ with $\pi_{is}(n_s^e) > 0$ and $i$
        was previously capacity-constrained, firm $i$ now enters $s$.
  \item Otherwise, firm $i$ redirects capacity to Santiago's
        above-threshold tenders.
\end{enumerate}

\bigskip
\textbf{Effect on Santiago (export regions):}\\
Redirected capacity increases entrants in above-threshold markets,
\emph{reducing} procurement costs there.

\medskip
\textbf{Effect on rural markets (import regions):}\\
Fewer non-local bidders in the 30--100 UTM band; competition falls,
prices rise. But local entry partially offsets this.

\bigskip
\small\textit{This cross-market linkage is the key contribution of the
multi-market framework --- single-market analyses miss it entirely.}
\end{frame}

% ── Welfare ──────────────────────────────────────────────
\begin{frame}{Model: Welfare}
\textbf{Government procurement surplus:}
\[
  W_r^{gov} = Q_r \cdot \left(\bar{p} - E[p_r]\right)
\]

\textbf{Local economic benefit:}
\[
  W_r^{local} = Q_r \cdot \Pr[\text{local wins in } r] \cdot w_r
\]
where $w_r$ captures local employment, wages, and multiplier effects.

\textbf{Total welfare:}
\[
  W = \sum_{r} \left[W_r^{gov} + \phi \cdot W_r^{local}\right]
\]
$\phi \geq 0$: social weight on local economic activity.

\textbf{Change from policy:}
\[
  \Delta W = \sum_r Q_r
    \left[
      \underbrace{-(E[\tilde{p}_r] - E[p_r])}_{\text{cost effect}}
      + \underbrace{\phi \cdot \Delta\!\Pr[\text{local wins}]_r \cdot w_r}_{\text{local benefit}}
    \right]
\]
\textbf{Sign is ambiguous and varies across regions} --- the core tension.
\end{frame}

% ── Implications vs. Empirics ─────────────────────────────
\begin{frame}{Model Implications vs.\ Current Empirical Results}
\begin{columns}[T]
  \column{0.5\textwidth}
  \textbf{Matched by the data:}
  \begin{itemize}
    \item[\checkmark] SME bidder share $\uparrow$,
          non-local share $\downarrow$ post-reform
    \item[\checkmark] SME winners $\uparrow$,
          large-firm winners $\downarrow$
    \item[\checkmark] Win prices $\uparrow$ (cost effect $>$ competition recovery)
    \item[\checkmark] Heterogeneity: regions with high pre-reform non-local
          share see largest local entry gains
    \item[\checkmark] Cross-market spillovers: above-threshold N bidders
          $\uparrow$ in high-export regions (capacity reallocation confirmed)
    \item[\checkmark] Large-firm spillover $>$ SME spillover
  \end{itemize}

  \column{0.5\textwidth}
  \textbf{Unresolved / needs more work:}
  \begin{itemize}
    \item[\textbf{?}] Entry response timing: how quickly do new local
          firms appear?
    \item[\textbf{?}] $\lambda$ vs.\ $\delta$: cannot yet separate
          distance cost from local informational advantage without
          structural estimation
    \item[\textbf{?}] $N_r^{pot}$ margin: limited data on potential
          (non-entering) local firms
    \item[\textbf{?}] Welfare weights $\phi$: distributional gains
          require employment data (SII match)
    \item[\textbf{?}] Bunching at thresholds: not yet tested
  \end{itemize}
\end{columns}
\end{frame}


% ══════════════════════════════════════════════════════════
% ESTIMATION PLAN
% ══════════════════════════════════════════════════════════
\section{Estimation Plan}

\begin{frame}{Estimation Plan: Overview}
\textbf{Goal:} Identify $(\delta, \lambda, \kappa_0, \kappa_d, K)$ using
variation in entry and bidding across regions, distances, and the reform.

\bigskip
\textbf{Three-step strategy:}
\begin{enumerate}
  \item \textbf{Reduced-form moments} (already estimated):\\
        Event-study and DiD coefficients provide model-consistent
        moments relating entry changes to pre-reform market structure.

  \vspace{0.3em}
  \item \textbf{Bidding stage} (first, conditional on entry):\\
        Estimate cost distribution $F(\cdot)$ from bid data using
        standard first-price auction methods (e.g., GPV 2000;
        Guerre, Perrigne, Vuong). Recovers $\delta$ and $\lambda$
        from the difference in cost distributions for local vs.\
        non-local firms at varying distances.

  \vspace{0.3em}
  \item \textbf{Entry stage} (second, using cost estimates):\\
        Back out entry costs $(\kappa_0, \kappa_d)$ and capacity $K$
        from the free-entry condition, using variation in $Q_r$,
        $N_r^{pot}$, and the reform as instruments.
\end{enumerate}
\end{frame}

\begin{frame}{Estimation Plan: Identifying Variation}
\begin{columns}[T]
  \column{0.5\textwidth}
  \textbf{Identifying $\delta$ (distance cost):}
  \begin{itemize}
    \item Within-tender variation: same project, different bidder origins
    \item Bid level $\sim$ distance of bidder's home region
    \item Pre-reform data; no confound from exclusion
    \item Controls: firm FE, tender FE, sector-region trends
  \end{itemize}

  \vspace{0.5em}
  \textbf{Identifying $\lambda$ (local advantage):}
  \begin{itemize}
    \item Discontinuity: otherwise identical firms, local vs.\ non-local
    \item Residual cost gap after controlling for $\delta \cdot d$
    \item Validation: should predict pre-reform win share differential
  \end{itemize}

  \column{0.5\textwidth}
  \textbf{Identifying $\kappa_0, \kappa_d$ (entry costs):}
  \begin{itemize}
    \item Cross-regional variation in entry rates vs.\ market size $Q_r$
    \item Free-entry condition pins down $\kappa_0$ given cost estimates
    \item $\kappa_d$: how entry by distant firms varies with distance
          (pre-reform)
  \end{itemize}

  \vspace{0.5em}
  \textbf{Identifying $K$ (capacity):}
  \begin{itemize}
    \item Post-reform spillovers: which firms enter more above-threshold
          tenders after exclusion from 30--100 UTM?
    \item Variation in export share $\times$ firm size
    \item Direct measure: change in number of markets per firm
  \end{itemize}
\end{columns}
\end{frame}

% ══════════════════════════════════════════════════════════
% COUNTERFACTUALS
% ══════════════════════════════════════════════════════════
\section{Counterfactuals}

\begin{frame}{Counterfactuals}
\textbf{[PLACEHOLDER --- describe planned counterfactuals]}

\medskip
Once structural parameters are estimated, three main counterfactuals:

\begin{enumerate}
  \item \textbf{Optimal local preference threshold by region}\\
        For each region $r$, find the threshold $\underline{v}_r^*$ that maximizes
        total welfare $W_r^{gov} + \phi \cdot W_r^{local}$.
        Quantify the cost of the uniform 100 UTM policy.

  \vspace{0.4em}
  \item \textbf{Welfare decomposition}\\
        Separate the cost effect (government pays more) from the
        local benefit effect (local firms win more) across regions.
        Map the spatial distribution of winners and losers from the policy.

  \vspace{0.4em}
  \item \textbf{Alternative policy designs}\\
        Bid preference $\alpha$ (weak form) vs.\ full exclusion (strong form).
        Which achieves more local wins at lower cost to the government?
        Optimal $\alpha$ by market thickness.
\end{enumerate}
\end{frame}


% ══════════════════════════════════════════════════════════
% NEXT STEPS
% ══════════════════════════════════════════════════════════
\section{Next Steps}

\begin{frame}{Next Steps}
\begin{columns}[T]
  \column{0.5\textwidth}
  \textbf{Empirics (short run)}
  \begin{itemize}
    \item Add descriptive facts about local/non-local firm bids
          (pre-reform patterns; geography of export activity)
    \item Fill summary statistics tables (procurement + SII)
    \item Extend event study window as more post-reform months arrive
    \item Test for bunching at 100 UTM threshold post-reform
    \item Track new local entrants (first-time bidders)
          by region post-reform
  \end{itemize}

  \column{0.5\textwidth}
  \textbf{Structural estimation}
  \begin{itemize}
    \item Match SII firm data to ChileCompra bidders (RUT linkage)
    \item Estimate bidding-stage cost model (GPV approach);
          recover $\delta, \lambda$ from local/non-local bid distributions
    \item Estimate entry model; recover $\kappa_0, \kappa_d$
          using pre-reform cross-regional variation
    \item Use post-reform spillovers to identify capacity $K$
    \item Validate: simulate model-predicted entry response and
          compare to event-study estimates
  \end{itemize}
\end{columns}

\bigskip
\begin{itemize}
  \item \textbf{Longer run:} obtain SII employment data for welfare
        decomposition; compute local multipliers $w_r$ by region
  \item \textbf{Writing:} working paper draft (target: Q3 2026)
\end{itemize}
\end{frame}


% ══════════════════════════════════════════════════════════
% APPENDIX
% ══════════════════════════════════════════════════════════

\begin{frame}
\centering
\textbf{APPENDIX}
\end{frame}

\appendix

% ── Monthly project counts and values ────────────────────

\begin{frame}{App: Monthly Project Counts by Band}
\begin{columns}[c]
  \column{0.49\textwidth}\centering
  \textbf{\small 0--30 UTM}\\[2pt]
  \includegraphics[width=\linewidth]{ca_D13_monthly_tender_counts_stacked_0_30utm.png}
  \column{0.49\textwidth}\centering
  \textbf{\small 30--100 UTM (Treated Band)}\\[2pt]
  \includegraphics[width=\linewidth]{ca_D13_monthly_tender_counts_stacked_30_100utm.png}
\end{columns}
\end{frame}

\begin{frame}{App: Monthly Project Counts --- Higher Bands}
\begin{columns}[c]
  \column{0.49\textwidth}\centering
  \textbf{\small 100--500 UTM}\\[2pt]
  \includegraphics[width=\linewidth]{ca_D13_monthly_tender_counts_stacked_100_500utm.png}
  \column{0.49\textwidth}\centering
  \textbf{\small 500+ UTM}\\[2pt]
  \includegraphics[width=\linewidth]{ca_D13_monthly_tender_counts_stacked_500_plus_utm.png}
\end{columns}
\end{frame}

\begin{frame}{App: Monthly Project Value by Band}
\begin{columns}[c]
  \column{0.49\textwidth}\centering
  \textbf{\small 0--30 UTM}\\[2pt]
  \includegraphics[width=\linewidth]{ca_D14_monthly_tender_value_stacked_0_30utm.png}
  \column{0.49\textwidth}\centering
  \textbf{\small 30--100 UTM (Treated Band)}\\[2pt]
  \includegraphics[width=\linewidth]{ca_D14_monthly_tender_value_stacked_30_100utm.png}
\end{columns}
\end{frame}

\begin{frame}{App: Monthly Project Value --- Higher Bands}
\begin{columns}[c]
  \column{0.49\textwidth}\centering
  \textbf{\small 100--500 UTM}\\[2pt]
  \includegraphics[width=\linewidth]{ca_D14_monthly_tender_value_stacked_100_500utm.png}
  \column{0.49\textwidth}\centering
  \textbf{\small 500+ UTM}\\[2pt]
  \includegraphics[width=\linewidth]{ca_D14_monthly_tender_value_stacked_500_plus_utm.png}
\end{columns}
\end{frame}

% ── Bidder types ─────────────────────────────────────────

\begin{frame}{App: Bidder Composition by Band}
\begin{columns}[c]
  \column{0.49\textwidth}\centering
  \textbf{\small 30--100 UTM}\\[2pt]
  \includegraphics[width=\linewidth]{ca_D8_bidder_types_avg_count_30_100utm_month.png}
  \column{0.49\textwidth}\centering
  \textbf{\small 100--500 UTM}\\[2pt]
  \includegraphics[width=\linewidth]{ca_D10_bidder_types_avg_count_100_500utm_month.png}
\end{columns}
\end{frame}

\begin{frame}{App: Bidder Share by Band (30--100 UTM)}
\centering
\includegraphics[width=0.90\linewidth]{ca_D20_bidder_type_share_30_100utm_month.png}
\end{frame}

\begin{frame}{App: Winner Share by Band (30--100 UTM)}
\centering
\includegraphics[width=0.90\linewidth]{ca_D24_winner_type_share_30_100utm_month.png}
\end{frame}

% ── OLS DiD Results ──────────────────────────────────────

\begin{frame}{App: OLS DiD --- Entry (Panel A, Municipalidades)}
\small\medskip
\input{../output/did/tables/did_table_compare_A_munic.tex}
\end{frame}

\begin{frame}{App: OLS DiD --- Bidder Composition (Panel B, Municipalidades)}
\small\medskip
\input{../output/did/tables/did_table_compare_B_munic.tex}
\end{frame}

\begin{frame}{App: OLS DiD --- Winner Characteristics (Panel C, Municipalidades)}
\small\medskip
\input{../output/did/tables/did_table_compare_C_munic.tex}
\end{frame}

\begin{frame}{App: OLS DiD --- Bid Outcomes (Panel D, Municipalidades)}
\small\medskip
\input{../output/did/tables/did_table_compare_D_munic.tex}
\end{frame}

% ── Heterogeneity details ─────────────────────────────────

\begin{frame}{App: Binscatter --- Market Thickness vs.\ Non-local Share}
\begin{center}
\includegraphics[width=0.90\linewidth]{binscatter_q_nonlocal_munic.png}
\end{center}
\small\textit{Thick markets (high q\_pre) have systematically fewer non-local
bidders --- log scale removes Santiago leverage.}
\end{frame}

\begin{frame}{App: Binscatter --- Moderator Pair Matrix}
\begin{center}
\includegraphics[width=0.85\linewidth]{binscatter_moderator_pairs_munic.png}
\end{center}
\small\textit{q\_pre and totval\_pre are highly correlated. n\_pot\_local is
a complementary channel. Santiago ($\star$) has extreme leverage on raw scale.}
\end{frame}

\begin{frame}{App: Interacted DiD --- Non-local Bidder Share}
\scriptsize\medskip
\resizebox{\textwidth}{!}{%
\input{../output/did/tables/hetero_interacted_nonlocal_share_pre_munic.tex}}
\end{frame}

\begin{frame}{App: Interacted DiD --- Potential Local Firms}
\small\medskip
\resizebox{\textwidth}{!}{%
\input{../output/did/tables/hetero_interacted_n_pot_local_munic.tex}}
\end{frame}

\begin{frame}{App: Interacted DiD --- Distance from Santiago}
\small\medskip
\resizebox{\textwidth}{!}{%
\input{../output/did/tables/hetero_interacted_dist_from_santiago_munic.tex}}
\end{frame}

% ── Option B robustness ───────────────────────────────────

\begin{frame}{App: Robustness --- Dec 2024 Excluded (Option B)}
\textbf{Motivation:} The reform took effect Dec 12. The $k=0$ month captures
only partial exposure ($\approx 19/31$ days treated).

\bigskip
\textbf{Option B:} Drop $k=0$ entirely. Post-period begins at $k=1$ (Jan 2025).
\end{frame}

\begin{frame}{App: Event Studies --- Option B}
\begin{columns}[c]
  \column{0.5\textwidth}\centering
  \includegraphics[width=\linewidth]{../output/did/figures/diag_event_study_n_bidders_optB.png}
  \column{0.5\textwidth}\centering
  \includegraphics[width=\linewidth]{../output/did/figures/diag_event_study_sme_share_sii_optB.png}
\end{columns}
\end{frame}

\begin{frame}{App: Pre-Trend Wald Test (Option B)}
\small\medskip
\input{../output/did/tables/diag_pretrend_wald_optB.tex}
\end{frame}

% ── Bid markups appendix ─────────────────────────────────
\begin{frame}{App: Bid Markups --- Residual Densities by Firm Type}
\centering
\includegraphics[width=0.96\linewidth]{bids_part1_resid_density_by_size.png}

\medskip
\small\textit{Within auction, local and non-local residual markups overlap heavily for SME and large firms. Non-SII bids are shown only in aggregate because locality is missing for unmatched firms.}
\end{frame}

\begin{frame}{App: Bid Markups --- Within-Firm Residual Density}
\centering
\includegraphics[width=0.82\linewidth]{bids_part2_resid_density_local_vs_nonlocal.png}

\medskip
\small\textit{After absorbing bidder and month effects, the residual distribution shifts left for local projects. The same firm tends to bid lower in its home region.}
\end{frame}

\begin{frame}{App: Bid Markups --- By Firm Home Region}
\centering
\includegraphics[width=0.92\linewidth]{bids_part2_region_split.png}

\medskip
\small\textit{Pooled origin-region interactions show that the local effect is heterogeneous across firm home regions, while the non-local distance slope is more consistently negative for larger central and southern origin regions. I read this as stronger evidence for a continuous distance friction than for one universal same-region premium.}
\end{frame}

\begin{frame}{App: Bid Markups --- Event Study}
\centering
\includegraphics[width=0.76\linewidth]{bids_part3_event_study.png}

\medskip
\small\textit{Post-reform bid markups jump sharply, but pre-period coefficients are not flat relative to $k=-1$. I treat this dynamic profile as descriptive rather than a clean event-study design.}
\end{frame}

\begin{frame}{App: Bid Markups --- Auction FE Specs}
\centering
\scriptsize
\resizebox{\linewidth}{!}{%
\input{../output/bids/tables/bids_part1_auction_fe.tex}
}
\end{frame}

\begin{frame}{App: Bid Markups --- Within-Firm and Distance Specs}
\centering
\scriptsize
\resizebox{\linewidth}{!}{%
\input{../output/bids/tables/bids_part2_firm_fe.tex}
}
\end{frame}

% ── Commune-level distance (new specs 3a–3f_com) ──────────
\begin{frame}{App: Commune-Level Distances --- Identification Motivation}
\begin{columns}[T]
  \column{0.5\textwidth}
  \textbf{Problem with region-centroid distances (specs 2d–2e):}
  \begin{itemize}
    \item Only 16 distinct values for non-local bids
    \item Spec (2e) bins are non-monotonic: 800+ km coefficient
          \emph{falls} relative to 400--800 km, likely due to strong
          positive selection among extreme-distance bidders
    \item $\delta$ (distance cost) and $\lambda$ (local info advantage)
          are confounded: both are identified off the same
          local/non-local boundary
  \end{itemize}

  \column{0.5\textwidth}
  \textbf{Solution: municipality (commune)-level distances:}
  \begin{itemize}
    \item $\sim$346 distinct values from full pairwise Haversine matrix
    \item Same-region bids have \emph{nonzero} commune distances
          $\Rightarrow$ within-region variation identifies $\delta$
          \emph{without} touching the local/non-local boundary
    \item Spec (3f\_com): local dummy + log distance on \emph{all} bids
          directly estimates the discrete local premium $\lambda$
          conditional on the continuous distance gradient $\delta$
    \item Near-border firms (same region, large commune distance)
          vs.\ same-region firms close to buyer sharpen the wedge
  \end{itemize}
\end{columns}

\medskip
\centering\small
\textit{New specs mirror 2a--2f but use commune centroids.
        Finer distance bins (1--50, 50--150, 150--300, 300--600, 600+ km)
        exploit the richer variation.}
\end{frame}

\begin{frame}{App: Commune-Level Distances --- Distance Gradient (Region vs.\ Commune)}
\centering
\includegraphics[width=0.86\linewidth]{bids_commune_distance_gradient.png}

\medskip
\small\textit{%
  Left: region-centroid log-distance slope on non-local bids (spec 2d).
  Right: commune-centroid log-distance slope on non-local bids (spec 3c\_com).
  Consistent slopes across both levels confirm the distance-cost gradient is not
  an artifact of the coarse regional grid.  Any attenuation in the commune-level
  estimate would indicate positive selection at long distances was inflating the
  region-level slope.%
}
\end{frame}

\begin{frame}{App: Commune-Level Distances --- Distance Bins (Region vs.\ Commune)}
\centering
\includegraphics[width=0.88\linewidth]{bids_commune_bins_comparison.png}

\medskip
\small\textit{%
  Left: region-centroid bins from spec (2e) --- non-monotonic at 800+ km.
  Right: finer commune-centroid bins from spec (3d\_com), reference = 1--50 km.
  If the 800+ km reversal disappears with finer bins, it was driven by extreme
  positive selection among very-long-distance bidders (consistent with
  $\delta > 0$ throughout); if it persists, it reflects a genuine structural
  non-monotonicity requiring a richer entry model.%
}
\end{frame}

\begin{frame}{App: Commune-Level Distances --- Regression Specs (3a--3f\_com)}
\centering
\scriptsize
\resizebox{\linewidth}{!}{%
\input{../output/bids/tables/bids_commune_distance.tex}
}

\medskip
\footnotesize\textit{%
  (3a\_com): local dummy + bidder \& month FE.
  (3b\_com): log commune distance on \emph{all} bids --- identifies $\delta$
  from within-region variation.
  (3c\_com): log commune distance, non-local only --- direct analogue of spec (2d).
  (3f\_com): local dummy + log commune distance on all bids --- direct joint
  estimate of $\lambda$ and $\delta$.
  (3d\_com): commune bins, non-local bids, ref = 1--50 km.
  (3e\_com): commune bins, all bids, ref = 0 (same commune).
  Clustered SE at bidder level throughout.
  Comparing (3b\_com) slope to (3c\_com) slope tests whether within-region
  and cross-region distance frictions are the same parameter $\delta$.%
}
\end{frame}

\begin{frame}{App: Commune-Level Distances --- Separating $\lambda$ from $\delta$}
\begin{columns}[T]
  \column{0.54\textwidth}
  \textbf{Reading specs (3a, 3b, 3c, 3f\_com) together:}
  \begin{enumerate}
    \item \textbf{Spec (3c\_com)} gives $\hat\delta$ from non-local bids,
          exactly as spec (2d) but at commune resolution.
          This is the clean cross-region benchmark.
    \item \textbf{Spec (3b\_com)} gives $\hat\delta_{\text{all}}$ from
          \emph{all} bids, including local within-region variation.
          If $\hat\delta_{\text{all}} \approx \hat\delta$, the
          distance friction is homogeneous and $\lambda$ is a
          pure informational premium on top of $\delta$.
    \item \textbf{Spec (3f\_com)} puts both in one regression:
          the local coefficient is the discrete same-region premium
          $\hat\lambda$ conditional on commune distance, while the
          slope is the continuous distance friction $\hat\delta$.
    \item \textbf{Comparing (3a\_com) to (3f\_com)} shows how much of the
          raw local gap was really shorter travel distance rather than a
          separate informational or relational premium.
  \end{enumerate}

  \column{0.46\textwidth}
  \textbf{What each result tells us structurally:}

  \smallskip
  \textit{If $\hat\delta_{\text{all}} \approx \hat\delta_{\text{nonloc}}$:}\\
  Distance cost is the same parameter whether we're inside or outside the
  region.  The local dummy absorbs only the informational piece $\lambda$.

  \smallskip
  \textit{If $\hat\delta_{\text{all}} < \hat\delta_{\text{nonloc}}$:}\\
  Local bids are less distance-sensitive even within the region $\Rightarrow$
  $\lambda$ partly reflects lower effective distance costs for local firms
  (e.g.\ existing supplier relationships closer to home).

  \smallskip
  \textit{If the local coefficient in (3f\_com) stays large:}\\
  There is a true discrete same-region advantage even after controlling for
  municipality-level distance.  If it collapses toward zero, the "local"
  effect was mostly a coarse-distance artifact.

  \smallskip
  \textit{Spec (3e\_com) bins on all bids:}\\
  Non-parametric check --- does the markup profile flatten inside the region
  (bin ``0'' vs.\ ``1--50'') or is it continuous?  A flat profile near zero
  indicates a discrete jump $\lambda$; a smooth profile indicates pure $\delta$.
\end{columns}
\end{frame}

\begin{frame}{App: Bid Markups --- Reform Specs}
\centering
\scriptsize
\resizebox{\linewidth}{!}{%
\input{../output/bids/tables/bids_part3_did.tex}
}
\end{frame}

% ── Model comparative statics ─────────────────────────────

\begin{frame}{App: When Does the Policy Reduce Procurement Costs?}
\begin{columns}[T]
  \column{0.5\textwidth}
  \textbf{Policy more likely to \emph{reduce} costs when:}
  \begin{enumerate}
    \item High $N_r^{pot}$ relative to $Q_r$: large local firm pool
          replaces excluded non-locals
    \item High $\lambda$: non-locals had local disadvantage anyway;
          their exclusion doesn't remove the best bidders
    \item High $\delta \cdot d_{sr}$: non-locals came from far away
          and were high-cost
    \item Low $K$ (tight capacity): freed capacity redirects home,
          generating competition spillovers
  \end{enumerate}

  \column{0.5\textwidth}
  \textbf{Policy more likely to \emph{increase} costs when:}
  \begin{enumerate}
    \item Low $N_r^{pot}$ relative to $Q_r$: few local firms,
          cannot replace competition
    \item Low $\lambda$, low $d_{sr}$: non-locals were genuinely
          competitive nearby firms
    \item High $K$ (ample capacity): no meaningful cross-market
          reallocation
    \item Lumpy demand ($Q_r$ volatile): local market too thin to
          sustain competition
  \end{enumerate}
\end{columns}
\end{frame}

\end{document}
